
\section{Compiler and Runtime Library}
\label{sec:compiler}

% Author voice. Figure shows graph node → ops.yaml → XCL emit.
% Long identifiers break at \allowbreak so they stay in-column.

\pyv{} starts from a \texttt{torch.export} graph of a {PyTorch} model.
Each node of the graph represents an operator, like a 2D convolution or a ReLU. 
While walking the graph, \pyv{} checks whether that operator is a lowering candidate,
and when it is, looks up the matching entry in an operation dictionary, containing
the mapping between PyTorch/ATEN operators and their custom {C} implementation in {XCL}.

That entry names a codegen handler and an {XCL} symbol (plus the {C} sources that implement it).
The handler emits a call into the generated {C} file, plus any pre- and post-processing needed.
After the codegen phase, the compiler hands it off to the {RISC-V} {GNU} toolchain, for final linking 
against {XCL} and {ELF} generation.

\begin{figure*}[t]
  \centering
  % Full-width lowering walk: graph → ops.yaml → generated call → XCL.
% Print-friendly: white ground, black ink, gray only for selection.
\definecolor{tvink}{HTML}{111111}
\definecolor{tvmute}{HTML}{333333}
\definecolor{tvline}{HTML}{555555}
\definecolor{tvshade}{HTML}{EEEEEE}

\usetikzlibrary{fit}
\pgfdeclarelayer{backcards}
\pgfdeclarelayer{backarrows}
\pgfsetlayers{backcards,backarrows,main}

\begin{tikzpicture}[
  >=Stealth,
  x=1mm, y=1mm,
  font=\sffamily\scriptsize,
  box/.style={
    draw=tvink,
    line width=0.6pt,
    rounded corners=1.5pt,
    fill=white,
    align=center,
    inner sep=2.4pt,
    font=\ttfamily\tiny,
  },
  leaf/.style={
    box,
    font=\ttfamily\tiny,
    minimum width=9mm,
  },
  selected/.style={
    box,
    line width=1.15pt,
    fill=tvshade,
  },
  card/.style={
    draw=tvink,
    line width=0.6pt,
    rounded corners=2pt,
    fill=white,
    align=left,
    inner xsep=3pt,
    inner ysep=2.6pt,
    font=\ttfamily\tiny,
  },
  hit/.style={
    font=\ttfamily\tiny,
    draw=tvink,
    line width=0.5pt,
    rounded corners=0.7pt,
    fill=tvshade,
    inner xsep=1.2pt,
    inner ysep=0.35pt,
  },
  arr/.style={
    ->,
    draw=tvline,
    line width=0.55pt,
    shorten >=1.0pt,
    shorten <=1.0pt,
  },
]
\pgfmathsetmacro{\W}{\textwidth/1mm - 2}

\draw[tvink, line width=0.6pt, rounded corners=2.5pt]
  (0,-2.2) rectangle (\W,51.5);

\node[font=\sffamily\footnotesize\bfseries, text=tvink, anchor=west]
  at (2.0,48.8) {One node through PyVedas};
\draw[tvink, line width=0.45pt]
  (2.0,46.6) -- (\W-2.0,46.6);

\node[font=\sffamily\footnotesize\bfseries, text=tvink, anchor=west]
  at (2.0,43.5) {1.\ Graph};
\node[font=\sffamily\footnotesize\bfseries, text=tvink, anchor=west]
  at (0.255*\W,43.5) {2.\ Select + look up};
\node[font=\sffamily\footnotesize\bfseries, text=tvink, anchor=west]
  at (0.50*\W,43.5) {3.\ Emit into C};
\node[font=\sffamily\footnotesize\bfseries, text=tvink, anchor=west]
  at (0.74*\W,43.5) {4.\ XCL function body};

\node[leaf] (n_x) at (0.04*\W, 34.0) {x};
\node[leaf] (n_w) at (0.19*\W, 34.0) {weight};
\node[box]  (n_conv) at (0.115*\W, 23.0) {conv2d};
\node[selected] (n_pool) at (0.115*\W, 12.0) {max\_pool2d};
\node[leaf] (n_y) at (0.115*\W, 2.2) {y};

\node[font=\sffamily\scriptsize\bfseries, text=tvink, anchor=south west]
  at ([xshift=1.2mm,yshift=0.6mm]n_pool.north east) {selected};

\node[font=\ttfamily\tiny, anchor=north west, align=left]
  (yl0) at (0.26*\W, 40.6) {\textcolor{tvmute}{\# ops.yaml}};
\node[font=\ttfamily\tiny, anchor=north west]
  (yl1) at ([yshift=-0.3mm]yl0.south west) {aten.max\_pool2d.default:};
\node[font=\ttfamily\tiny, anchor=north west]
  (yl2) at ([yshift=-0.3mm]yl1.south west) {\quad codegen: max\_pool2d};
\node[font=\ttfamily\tiny, anchor=north west]
  (yl3) at ([yshift=-0.3mm]yl2.south west) {\quad symbol:};
\node[hit, anchor=north west]
  (sym) at ([xshift=4.0mm,yshift=-0.4mm]yl3.south west)
  {pyvedas\_aten\_max\_pool2d};
\node[font=\ttfamily\tiny, anchor=north west]
  (yl4) at ([yshift=-1.0mm]yl0.west |- sym.south) {\quad sources:};
\node[font=\ttfamily\tiny, anchor=north west]
  (yl5) at ([yshift=-0.3mm]yl4.south west) {\qquad aten\_max\_pool2d.c};

\node[font=\ttfamily\tiny, anchor=north west, align=left]
  (cl0) at (0.505*\W, 40.6) {\textcolor{tvmute}{/* generated */}};
\node[hit, anchor=north west]
  (call) at ([yshift=-0.3mm]cl0.south west) {pyvedas\_aten\_max\_pool2d(};
\node[font=\ttfamily\tiny, anchor=north west]
  (cl2) at ([yshift=-0.3mm]call.south west) {\quad x, out, n, c, h, w,};
\node[font=\ttfamily\tiny, anchor=north west]
  (cl3) at ([yshift=-0.3mm]cl2.south west) {\quad oh, ow, kh, kw,};
\node[font=\ttfamily\tiny, anchor=north west]
  (cl4) at ([yshift=-0.3mm]cl3.south west) {\quad sh, sw, pad\_h, pad\_w);};

\node[font=\ttfamily\tiny, anchor=north west, align=left]
  (xl0) at (0.745*\W, 40.6) {\textcolor{tvmute}{/* aten\_max\_pool2d.c */}};
\node[font=\ttfamily\tiny, anchor=north west]
  (xclfn) at ([yshift=-0.3mm]xl0.south west) {void pyvedas\_aten\_max\_pool2d(};
\node[font=\ttfamily\tiny, anchor=north west]
  (xl2) at ([yshift=-0.3mm]xclfn.south west) {\quad const int32\_t *x,};
\node[font=\ttfamily\tiny, anchor=north west]
  (xl3) at ([yshift=-0.3mm]xl2.south west) {\quad int32\_t *out, \ldots)};
\node[font=\ttfamily\tiny, anchor=north west]
  (xl4) at ([yshift=-0.3mm]xl3.south west) {\{};
\node[font=\ttfamily\tiny, text=tvmute, anchor=north west]
  (xl5) at ([yshift=-0.3mm]xl4.south west) {\quad /* max-pool body here */};
\node[font=\ttfamily\tiny, anchor=north west]
  (xl6) at ([yshift=-0.3mm]xl5.south west) {\}};

\begin{pgfonlayer}{backcards}
  \node[card, fit=(yl0)(yl1)(yl2)(yl3)(sym)(yl4)(yl5),
    inner xsep=3.2pt, inner ysep=2.4pt] (yaml) {};
  \node[card, fit=(cl0)(call)(cl2)(cl3)(cl4),
    inner xsep=3.2pt, inner ysep=2.4pt] (cout) {};
  \node[card, fit=(xl0)(xclfn)(xl2)(xl3)(xl4)(xl5)(xl6),
    inner xsep=3.2pt, inner ysep=2.4pt] (xcl) {};
\end{pgfonlayer}

\begin{pgfonlayer}{backarrows}
  \draw[arr] (n_x.south) -- (n_conv.north west);
  \draw[arr] (n_w.south) -- (n_conv.north east);
  \draw[arr] (n_conv.south) -- (n_pool.north);
  \draw[arr] (n_pool.south) -- (n_y.north);
  \draw[arr, line width=0.6pt]
    (n_pool.east) -- ([xshift=-4.5mm]yaml.west |- n_pool.east)
    -- ([xshift=-4.5mm,yshift=1.2mm]yaml.west)
    -- ([yshift=1.2mm]yaml.west);
  \draw[arr, line width=0.6pt]
    (sym.east) -- (call.west);
  \draw[arr, line width=0.6pt]
    (call.east) -- (xclfn.west);
\end{pgfonlayer}

\end{tikzpicture}

  \caption{Lowering one graph node.
    \pyv{} selects \texttt{aten.max\_pool2d.default} from a small
    exported graph, looks up the {XCL} candidate in \texttt{ops.yaml},
    and emits the call that the stock toolchain will compile.}
  \label{fig:pyvedas-lower}
\end{figure*}

Figure~\ref{fig:pyvedas-lower} shows what a graph walk looks like:
\texttt{conv2d} then \texttt{max\_pool2d}.
The highlighted pool node resolves to an {XCL} symbol in
\texttt{aten\_max\_pool2d.c}.
Convolution follows the same path, except its {XCL} body can call a
{DA} on the {XDAI} (the $8{\times}8$ {GEMM} in this paper) instead of
letting GCC unroll the operator to plain RISC-V instructions.
A better {XCL} body---or one that calls an available {DA}---is the
same lookup; \pyv{} still has to emit the call.

\pyv{} extracts weights from the graph by looking for \texttt{get\_attr}
nodes and materializes them into the compile image (static buffers or the
on-chip memory), so the generated Executable and Linkable Format
({ELF}) carries both code and parameters.\footnote{The careful reader will notice how there is nothing special about RISC-V.
An ARM licensee can easily swap RISC-V GCC for ARM GCC, and swap the RISC-V core
with their ARM core and the system would work the same. C is the intermediate
representation language.}


